\begin{proposition}[6.11.8]
\label{IV.2.s6.11.8-fr}
Sea $X$ un preesquema localmente noetheriano que verifica la propiedad siguiente:
\begin{description}
\item[\textmd{(CMU)}] Todo subpreesquema cerrado íntegro $Y$ de $X$ contiene un abierto
no vacío $W$ tal que el preesquema inducido por $Y$ sobre $W$ es un preesquema
de Cohen-Macaulay.
\end{description}
Entonces, para todo $\mathcal O_X$-Módulo coherente $\mathcal F$, la función
$x\leadsto\operatorname{coprof}(\mathcal F_x)$ es localmente construible
y semicontinua superiormente; los conjuntos $U_{C_n}(\mathcal F)$ y
$U_{S_n}(\mathcal F)$ son abiertos en $X$.
\end{proposition}

\begin{proof}
En efecto, sea $Y$ un subpreesquema cerrado íntegro de $X$ con punto
genérico $y$; en virtud de (6.10.6), existe en $Y$ una vecindad abierta
$V$ de $y$ tal que, para todo $x\in V\cap Y$, se tiene
\begin{equation}
\operatorname{coprof}(\mathcal F_x)=
\operatorname{coprof}(\mathcal F_y)+
\operatorname{coprof}(\mathcal O_{Y,x}).
\tag{6.11.8.1}\label{IV.2.s6.11.8.1-fr}
\end{equation}
Pero, por hipótesis, existe un abierto no vacío $W$ de $Y$ tal que, para
$x\in V\cap W$, se tiene $\operatorname{coprof}(\mathcal O_{Y,x})=0$; por tanto,
$\operatorname{coprof}(\mathcal F_x)$ es constante en una vecindad de
$y$ en $Y$, lo que prueba que la función
$x\leadsto\operatorname{coprof}(\mathcal F_x)$ es localmente construible
(\textbf{0}$_{\rm III}$, 9.3.2). Además, se deduce de (6.11.5) y de
(\textbf{0}$_{\rm III}$, 9.3.4) que esta función es semicontinua superiormente.
La última afirmación se sigue de (6.11.4, (i)), o también de (6.11.6).
\end{proof}

\begin{remarks}[6.11.9]
\label{IV.2.s6.11.9-fr}
\begin{enumerate}
\item[(i)] Si $X$ verifica la hipótesis (CMU) de (6.11.8), entonces, para todo
morfismo $u:X'\to X$ \emph{localmente de tipo finito}, $X'$ también verifica (CMU).
Sean, en efecto, $Y'$ un subpreesquema cerrado íntegro de $X'$, $y'$ su punto
genérico, $y=u(y')$, y sea $Y$ el subpreesquema íntegro de $X$ cuyo
espacio subyacente es $\overline{\{y\}}$; entonces $u|Y'$ se factoriza como
$Y'\xrightarrow{v}Y\xrightarrow{j}X$, donde $j$ es la inyección canónica
(\textbf I, 5.2.2), y $v$ es localmente de tipo finito (\textbf I, 6.6.6). Al restringirse
a vecindades abiertas afines de $y$ y $y'$ respectivamente,
es posible, por tanto, limitarse al caso en que $X=\operatorname{Spec}(A)$, donde $A$ es un
anillo íntegro de Cohen-Macaulay, y $X'=\operatorname{Spec}(A')$, donde $A'$
es un anillo
\oldpage[IV-2]{163}
íntegro que contiene a $A$ y es una $A$-álgebra de tipo finito.
Reemplazando, si es necesario, $A$ por un anillo de fracciones $A_f$ (con $f\ne0$),
se puede suponer, además, que $A'$ contiene un anillo de polinomios
$A[T_1,\ldots,T_n]=A''$ y es una $A''$-álgebra \emph{finita} (Bourbaki,
\emph{Alg. comm.}, cap. V, \S3, n\textsuperscript{o} 1, cor. 1 del teorema 1).
Pero $A''$ es un anillo de Cohen-Macaulay (6.3.6); por tanto, es posible limitarse
al caso en que, además, $A'$ es una $A$-álgebra \emph{finita}. Existe entonces $g\ne0$ en
$A$ tal que $A'_g$ es un $A_g$-módulo \emph{libre} finitamente generado (Bourbaki,
\emph{Alg. comm.}, cap. II, \S5, n\textsuperscript{o} 1, cor. de la
prop. 2); por tanto, se puede suponer también que $A'$ es un $A$-módulo
\emph{libre}. Pero entonces $A'$ es un $A$-módulo de Cohen-Macaulay (\textbf 0, 16.5.1), y
como $A'$ es un $A$-módulo finitamente generado, $A'$ es también un $A'$-módulo de
Cohen-Macaulay (\textbf 0, 16.5.3), es decir, un anillo de Cohen-Macaulay.
\item[(ii)] Supongamos que existe un $\mathcal O_X$-Módulo coherente
$\mathcal F$ tal que $\operatorname{Supp}(\mathcal F)=X$ y que
$\mathcal F$ es un $\mathcal O_X$-Módulo \emph{de Cohen-Macaulay}. Entonces $X$
verifica la condición (CMU): en efecto, con las notaciones de la demostración
de (6.11.8), la relación (6.11.8.1) muestra que se tiene
$\operatorname{coprof}(\mathcal O_{Y,x})=0$ en una vecindad (relativa a
$Y$) del punto genérico $y$ de $Y$.
\end{enumerate}

Se ignora si existen preesquemas localmente noetherianos de dimensión
$\geq2$ que no verifiquen (CMU) (si $\dim(X)=1$, es inmediato que todo
punto maximal de $X_{\mathrm{red}}$ admite una vecindad abierta íntegra de
dimensión $1$ y, por tanto, de Cohen-Macaulay).
\end{remarks}
